% Questions et corrigé du Jour 4, Série 2 (Automatismes Hebdo). Document
% autonome, distinct du livret LaTeX complet et de la "Feuille élève -
% Série 2 - Semaine du 14 au 18 septembre.tex". Même gabarit que
% "Série2-J1/J2/J3 - Questions et corrigé.tex". Se compile seul avec
% pdflatex ou lualatex.
\documentclass[12pt]{article}
\usepackage[a4paper,landscape,top=6mm,bottom=6mm,left=10mm,right=10mm]{geometry}
\usepackage[T1]{fontenc}
\usepackage[french]{babel}
\usepackage[table]{xcolor}
\usepackage[most]{tcolorbox}
\usepackage{ragged2e}
\usepackage{amsmath}
\usepackage{pifont}
\usepackage{enumitem}
\usepackage{tikz}

\setlength{\parindent}{0pt}
\renewcommand{\familydefault}{\sfdefault}

% --- Palette (reprise de "Feuille élève - Série 2 - Semaine du 14 au 18 septembre.tex") ---
\definecolor{inkC}{HTML}{20232B}
\definecolor{inksoftC}{HTML}{565A66}
\definecolor{mutedC}{HTML}{8B8F99}
\definecolor{qcmC}{HTML}{2455B8}
\definecolor{questionC}{HTML}{12805C}
\definecolor{finalC}{HTML}{B5511F}
\definecolor{paperC}{HTML}{F2F3EE}
\definecolor{consoleBgC}{HTML}{1E1E2E}
\definecolor{consoleGreenC}{HTML}{6FCF97}
\definecolor{tacheBlueC}{HTML}{2D9CDB}

\newcommand{\okmark}{\ding{51}}

% Console façon terminal (fond sombre, invite ">>>" en vert) : le résultat
% reste visible (document projeté en classe), simplement surligné d'une
% "tache" bleue façon surligneur, pour attirer l'œil dessus.
\newcommand{\consoleBoxResult}[2]{%
  % #1 = appel affiché après >>>, #2 = résultat surligné
  \begin{tcolorbox}[colback=consoleBgC, colframe=consoleBgC, arc=3pt,
    boxrule=0pt, left=8pt, right=8pt, top=3pt, bottom=3pt]
  \ttfamily\footnotesize\color{white}
  {\color{consoleGreenC}>>>} \# Script executed\\[2pt]
  {\color{consoleGreenC}>>>} #1\\[5pt]
  \begin{tikzpicture}[baseline=-2pt]
    \fill[tacheBlueC, opacity=0.6]
      (0,0)
      .. controls (0.2,0.24) and (0.55,-0.1) .. (0.9,0.11)
      .. controls (1.2,0.28) and (1.5,-0.07) .. (1.78,0.12)
      .. controls (1.98,0.25) and (1.78,0.42) .. (1.42,0.36)
      .. controls (1.05,0.46) and (0.45,0.44) .. (0.15,0.28)
      .. controls (-0.06,0.19) and (-0.06,0.07) .. (0,0) -- cycle;
    \node[white, font=\bfseries] at (0.9,0.16) {#2};
  \end{tikzpicture}
  \end{tcolorbox}%
}

\NewDocumentCommand{\RectoTitre}{m m}{%
  \begin{center}
  \begin{tcolorbox}[
    enhanced, width=0.85\linewidth, colback=white, colframe=inkC,
    boxrule=0.8pt, arc=8pt, sharp corners=downhill,
    left=10pt, right=10pt, top=4pt, bottom=3pt, halign=center,
    overlay={
      \node[fill=white, anchor=west, font=\normalsize, inner xsep=4pt]
        at ([xshift=14pt]frame.north west) {#1};
    },
  ]
  {\huge\bfseries #2}
  \end{tcolorbox}
  \end{center}
}

\newcommand{\QuestionCorrigee}[5]{%
  \begin{tcolorbox}[
    enhanced, colback=white, colframe=#2, boxrule=1pt, arc=4pt,
    left=10pt, right=10pt, top=3pt, bottom=3pt,
    before skip=3pt, after skip=3pt,
  ]
  \noindent{\textcolor{#2}{\textbf{Question #1 \textperiodcentered{} #3}}}\hfill{\small\textcolor{inksoftC}{#4}}\par
  \vspace{3pt}
  #5
  \end{tcolorbox}%
}

% Petit bloc de code façon "console" du site (fond clair, cadre, monospace).
\newcommand{\codebox}[1]{%
  \fcolorbox{mutedC}{paperC}{\begin{minipage}{9.5cm}\ttfamily\footnotesize #1\end{minipage}}%
}

\begin{document}
\pagestyle{empty}

\normalsize
\RectoTitre{Série 2 --- Jour 4 --- Semaine du 14 au 18 septembre}{Questions et corrigé}
\vspace{4pt}

\begin{tcbraster}[raster columns=2, raster column skip=10pt, raster row skip=3pt, raster force size=false, raster valign=top]
\QuestionCorrigee{1}{qcmC}{QCM}{Vecteurs \textperiodcentered{} Égalité}{%
  On sait que $\overrightarrow{\text{AB}}=\overrightarrow{\text{DC}}$. Quel quadrilatère est un parallélogramme~?
  \vspace{3pt}
  \begin{itemize}[label=,leftmargin=1.2em,itemsep=1pt]
    \item[\textcolor{questionC}{\okmark}~A.] ABCD \hfill\textit{\small(réponse correcte)}
    \item[B.] ABDC
    \item[C.] ADCB
    \item[D.] ACBD
  \end{itemize}
  \vspace{1pt}
  \textbf{\textcolor{qcmC}{Corrigé}} --- Dire que $\overrightarrow{\text{AB}}=\overrightarrow{\text{DC}}$ signifie que $[\text{AB}]$ et $[\text{DC}]$ sont parallèles, de même longueur et de même sens~: ce sont deux côtés opposés d'un parallélogramme. En respectant l'ordre des sommets, c'est le quadrilatère \textbf{ABCD} qui est un parallélogramme.
}
\QuestionCorrigee{2}{qcmC}{QCM}{A3 --- Puissances}{%
  Simplifier $x^{3}\times x^{5}$ (le résultat sera écrit sous la forme $x^n$).
  \vspace{3pt}
  \begin{itemize}[label=,leftmargin=1.2em,itemsep=1pt]
    \item[A.] $x^{15}$
    \item[\textcolor{questionC}{\okmark}~B.] $x^{8}$ \hfill\textit{\small(réponse correcte)}
    \item[C.] $x^{2}$
    \item[D.] $2x^{8}$
  \end{itemize}
  \vspace{1pt}
  \textbf{\textcolor{qcmC}{Corrigé}} --- Pour multiplier deux puissances de même base, on additionne les exposants~: $x^3\times x^5=x^{3+5}=x^8$. $x^{15}$ viendrait d'une multiplication des exposants (confusion avec $(x^3)^5$)~; $x^2$ d'une soustraction (confusion avec la règle du quotient)~; $2x^8$ ajoute un coefficient qui n'a pas lieu d'être ici.
}
\QuestionCorrigee{3}{finalC}{Automatisme final}{A4 --- Écritures d'un nombre}{%
  Écrire en notation scientifique\\
  $2\times10^{3}+5\times10^{1}+3\times10^{-2}+8\times10^{3}$.
  \vspace{3pt}

  \textbf{\textcolor{finalC}{Corrigé}} --- On regroupe d'abord les termes en $10^{3}$~: $2\times10^{3}+8\times10^{3}=10\times10^{3}=10\,000$. Puis $5\times10^{1}=50$ et $3\times10^{-2}=0,03$. La somme vaut $10\,000+50+0,03=10\,050,03$, soit $1,005003\times10^{4}$ en notation scientifique.
}
\QuestionCorrigee{4}{questionC}{Question}{Calcul littéral \textperiodcentered{} Distributivité}{%
  Développer et réduire $4(1-x)+3x$.
  \vspace{3pt}

  \textbf{\textcolor{questionC}{Corrigé}} --- $4(1-x)+3x=4-4x+3x=4-x$.
}
\QuestionCorrigee{5}{questionC}{Question}{Algo \textperiodcentered{} Lire un programme}{%
  Voici un programme Python. Qu'affiche la console\\
  \textcolor{orange}{si on saisit \texttt{mystere(3, 9)}~?}
  \vspace{4pt}
  \begin{center}
  \fcolorbox{mutedC}{paperC}{\begin{minipage}{11.5cm}\ttfamily\small def mystere(a, b):\\
  \hspace*{1em}return a**2 + 2*b + 6\end{minipage}}

  \vspace{4pt}
  \consoleBoxResult{mystere(3, 9)}{33}
  \end{center}
  \vspace{1pt}
  \textbf{\textcolor{questionC}{Corrigé}} --- On remplace $a$ par $3$ et $b$ par $9$~: $a^2+2b+6=3^2+2\times9+6=9+18+6=33$.
}
\end{tcbraster}

\end{document}
